\begin{exo}{}
	Trois étudiants, Alice, Bob et Chloé, doivent se retrouver à la bibliothèque universitaire pour travailler sur un projet de mathématiques.

	La bibliothèque est ouverte en continu de $8h$ à $18h$.

	Chaque étudiant a des disponibilités bien précises au cours de cette journée :
	\begin{itemize}
		\item Alice est disponible entre $9h$ et $14h$.
		\item Bob est disponible à partir de $11h$ jusqu'à $16h$.
		\item Chloé est disponible soit entre $8h30$ et $10h$, soit entre $13h$ et $17h$.
	\end{itemize}
	\begin{enumerate}
		\item Écrirer les plages horaires de disponibilité d'Alice, de Bob et de Chloé sous la forme d'intervalles ou de réunions d'intervalles, que l'on notera respectivement $A$, $B$ et $C$.
		\item Représenter ces trois ensembles sur un même axe gradué en utilisant trois couleurs différentes.
		\item Alice et Bob décident de commencer à travailler ensemble en attendant Chloé.
		      \begin{enumerate}
			      \item Traduire mathématiquement leur plage horaire commune à l'aide d'un symbole d'opération sur les ensembles.
			      \item Déterminer cet ensemble sous forme d'intervalle, puis le traduire par une phrase claire en donnant les horaires correspondants.
		      \end{enumerate}
		\item Déterminer l'ensemble $A \cap B \cap C$. Que représente cet ensemble pour les trois étudiants ? Est-il possible qu'ils se réunissent tous les trois en même temps ?
		\item Déterminer l'ensemble $A \cup B$. À quoi correspond cet ensemble dans le contexte de l'exercice ?
	\end{enumerate}
\end{exo}
